\section{Vorhandene UX-Strategien und Gestaltungsprinzipien}
\label{sec:03-05_existing-ux-strategies}

Zum Zeitpunkt der Voranalyse existierten keine explizit dokumentierten \Gls{uxstrategien} oder verbindlichen Gestaltungsrichtlinien für die Plattform. Entscheidungen zu Struktur, Navigation und Darstellung erfolgten überwiegend dezentral und reaktiv; ein übergreifendes Rahmenkonzept für \Gls{ia}, Seitentypen, Navigationseinstiege oder visuelle Gestaltung war nicht etabliert, obwohl Forschung klare Design-Guidelines und zentrale UX-Governance als wesentliche Erfolgsfaktoren konsistenter Nutzungserlebnisse beschreibt \cite{113:NNGIntranetGuidelines,115:NNGContentManagement}.

Zwar waren einzelne \Gls{ux}-nahe Elemente implizit vorhanden, etwa wiederkehrende Seitentypen, visuelle Hervorhebungen oder strukturierte Einstiegspunkte, diese blieben jedoch punktuell und isoliert. Ohne gemeinsames Design-System sowie klar definierte Rollen und Verantwortlichkeiten für \Gls{ux} und \Gls{governance} fehlte eine systematische Verzahnung von Struktur, Gestaltung und inhaltlicher Pflege, wodurch das Verbesserungspotenzial dieser Ansätze nur eingeschränkt wirksam wurde \cite{113:NNGIntranetGuidelines,116:StepTwoGuidance,309:NNGContentGovernance}.

Diese Befunde werden durch ein im Praxisprojekt geführtes Experteninterview mit Prof. Dr. Peter Riegler gestützt. Er hob hervor, dass fehlende konsistente Strukturen, unklare Navigation und eine unzureichende visuelle Gestaltung zentrale Akzeptanzhemmnisse von Wikis darstellen, insbesondere für weniger technikaffine Nutzergruppen. Solche UX-Defizite verstärken psychologische Barrieren, etwa Unsicherheiten hinsichtlich der Bearbeitbarkeit von Inhalten, während klar definierte Seitentypen, kontextbezogene Verlinkungen und ein konsistentes visuelles Erscheinungsbild Orientierung, Vertrauen und wahrgenommene Nutzbarkeit fördern. Insgesamt unterstreichen diese Einschätzungen die Notwendigkeit einer expliziten UX-Strategie und verbindlicher Gestaltungsprinzipien.
